\sectionheading{5.33 Python Variable}

Program-এ Data বা Value-কে একটি নাম দিয়ে সংরক্ষণ ও ব্যবহার
করার জন্য \textbf{Variable} ব্যবহৃত হয়।

Python-এ C/C++-এর মতো Variable তৈরির আগে আলাদাভাবে
Data Type ঘোষণা করতে হয় না।

%================================================
% Definition
%================================================

\booksubheading{Python Variable কী?}

\begin{definitionbox}

\textbf{\color{darkTeal}সংজ্ঞা}

Python Program-এ কোনো Value-কে একটি নির্দিষ্ট নামের মাধ্যমে
সংরক্ষণ ও ব্যবহার করার জন্য যে নাম ব্যবহার করা হয় তাকে
\textbf{Variable} বলে।

\end{definitionbox}

উদাহরণ—

\begin{codebox}

\begin{english}
\ttfamily
age = 18
\end{english}

\end{codebox}

এখানে—

\begin{center}

\textenglish{\ttfamily age}
\(\longrightarrow\)
\textbf{Variable Name}
\qquad
\textenglish{\ttfamily 18}
\(\longrightarrow\)
\textbf{Value}
\qquad
\textenglish{\ttfamily =}
\(\longrightarrow\)
\textbf{Assignment Operator}

\end{center}

%================================================
% Creating Variable
%================================================

\booksubheading{Python-এ Variable তৈরি}

Python-এ কোনো নামের সঙ্গে Value Assign করলেই Variable তৈরি হয়।

উদাহরণ—

\begin{codebox}

\begin{english}
\ttfamily
x = 10\\
name = "Rahim"
\end{english}

\end{codebox}

এখানে
\textenglish{\ttfamily x}
এবং
\textenglish{\ttfamily name}
দুটি Variable।

\begin{rememberbox}

\textbf{\color{deepOrange}মনে রাখবে}

Python-এ সাধারণভাবে—

\begin{codebox}

\begin{english}
\ttfamily
int x
\end{english}

\end{codebox}

এভাবে Variable Declare করতে হয় না।

বরং—

\begin{codebox}

\begin{english}
\ttfamily
x = 10
\end{english}

\end{codebox}

লিখলেই Variable তৈরি হয়।

\end{rememberbox}

%================================================
% Assignment
%================================================

\booksubheading{Assignment}

Python-এ
\textenglish{\ttfamily =}
চিহ্ন ব্যবহার করে Variable-এর সঙ্গে Value Assign করা হয়।

উদাহরণ—

\begin{codebox}

\begin{english}
\ttfamily
number = 50
\end{english}

\end{codebox}

Assignment-এর সাধারণ রূপ—

\begin{center}

\textbf{Variable}
\(\;=\;\)
\textbf{Value}

\end{center}

অর্থাৎ—

\begin{center}

\textbf{ডান পাশের Value}
\(\longrightarrow\)
\textbf{বাম পাশের Variable}

\end{center}

উদাহরণ—

\begin{codebox}

\begin{english}
\ttfamily
price = 500
\end{english}

\end{codebox}

এখানে
\textenglish{\ttfamily 500}
Value-টি
\textenglish{\ttfamily price}
Variable-এর সঙ্গে Assign করা হয়েছে।

%================================================
% Using Variable
%================================================

\booksubheading{Variable-এর Value ব্যবহার}

Variable তৈরি করার পর Variable-এর নাম ব্যবহার করে তার
Value নিয়ে কাজ করা যায়।

\begin{examplebox}{উদাহরণ}

\begin{codebox}

\begin{english}
\ttfamily
age = 18\\
print(age)
\end{english}

\end{codebox}

\textbf{Output:}

\begin{codebox}

\begin{english}
\ttfamily
18
\end{english}

\end{codebox}

\end{examplebox}

%================================================
% Changing Value
%================================================

\booksubheading{Variable-এর Value পরিবর্তন}

একই Variable-এ পরে নতুন Value Assign করা যায়।

\begin{codebox}

\begin{english}
\ttfamily
x = 10\\
x = 20\\
print(x)
\end{english}

\end{codebox}

\textbf{Output:}

\begin{codebox}

\begin{english}
\ttfamily
20
\end{english}

\end{codebox}

শেষের Assignment-এ
\textenglish{\ttfamily x}
-এর নতুন Value হয়েছে
\textenglish{\ttfamily 20}।

\begin{center}

\textenglish{\ttfamily x = 10}
\(\longrightarrow\)
\textenglish{\ttfamily x = 20}
\(\longrightarrow\)
\textbf{বর্তমান Value = 20}

\end{center}

%================================================
% Naming Rules
%================================================

\booksubheading{Python Variable নামকরণের নিয়ম}

Variable-এর নাম লেখার ক্ষেত্রে কিছু নির্দিষ্ট নিয়ম মানতে হয়।

%------------------------------------------------
% Rule 1
%------------------------------------------------

\textbf{\color{darkTeal}১. Letter বা Underscore দিয়ে শুরু করা যায়}

\textbf{সঠিক—}

\begin{codebox}

\begin{english}
\ttfamily
name = "Rafi"\\
\_age = 18
\end{english}

\end{codebox}

\textbf{ভুল—}

\begin{codebox}

\begin{english}
\ttfamily
2name = "Rafi"
\end{english}

\end{codebox}

কারণ Variable Name Digit দিয়ে শুরু হতে পারে না।

%------------------------------------------------
% Rule 2
%------------------------------------------------

\textbf{\color{darkTeal}২. Letter, Digit ও Underscore ব্যবহার করা যায়}

সঠিক উদাহরণ—

\begin{codebox}

\begin{english}
\ttfamily
student1 = "Rahim"\\
roll\_10 = 25\\
total\_marks = 450
\end{english}

\end{codebox}

Variable Name-এ Letter, Digit ও Underscore ব্যবহার করা যায়;
তবে Digit প্রথম Character হতে পারে না।

%------------------------------------------------
% Rule 3
%------------------------------------------------

\textbf{\color{darkTeal}৩. Space ব্যবহার করা যায় না}

\textbf{ভুল—}

\begin{codebox}

\begin{english}
\ttfamily
student name = "Rahim"
\end{english}

\end{codebox}

\textbf{সঠিক—}

\begin{codebox}

\begin{english}
\ttfamily
student\_name = "Rahim"
\end{english}

\end{codebox}

\begin{rememberbox}

\textbf{\color{deepOrange}মনে রাখবে}

\begin{center}

\textbf{Space নয়}
\(\longrightarrow\)
\textbf{Underscore \textenglish{\ttfamily \_} ব্যবহার করা যায়}

\end{center}

\end{rememberbox}

%------------------------------------------------
% Rule 4
%------------------------------------------------

\textbf{\color{darkTeal}৪. Special Character ব্যবহার করা যায় না}

Variable Name-এ সাধারণভাবে নিচের ধরনের Special Character
ব্যবহার করা যায় না—

\begin{center}

\textenglish{\ttfamily @ \quad \# \quad \$ \quad \% \quad -}

\end{center}

\textbf{ভুল—}

\begin{codebox}

\begin{english}
\ttfamily
total-marks = 500\\
price\$ = 100
\end{english}

\end{codebox}

\textbf{সঠিক—}

\begin{codebox}

\begin{english}
\ttfamily
total\_marks = 500\\
price = 100
\end{english}

\end{codebox}

%------------------------------------------------
% Rule 5
%------------------------------------------------

\textbf{\color{darkTeal}৫. Python Keyword Variable Name হতে পারে না}

Python-এর Reserved Keyword Variable-এর নাম হিসেবে ব্যবহার
করা যায় না।

\textbf{ভুল—}

\begin{codebox}

\begin{english}
\ttfamily
if = 10\\
for = 20
\end{english}

\end{codebox}

কারণ
\textenglish{\ttfamily if}
ও
\textenglish{\ttfamily for}
Python Keyword।

\begin{rememberbox}

\textbf{\color{deepOrange}MCQ তথ্য}

\begin{center}

\textbf{Python Keyword}
\(\longrightarrow\)
\textbf{Variable Name হিসেবে ব্যবহার করা যায় না}

\end{center}

\end{rememberbox}

%------------------------------------------------
% Rule 6
%------------------------------------------------

\textbf{\color{darkTeal}৬. Python Variable Case Sensitive}

Python Variable Name Case Sensitive।

\begin{codebox}

\begin{english}
\ttfamily
age = 18\\
Age = 20\\
AGE = 25
\end{english}

\end{codebox}

এখানে
\textenglish{\ttfamily age},
\textenglish{\ttfamily Age}
এবং
\textenglish{\ttfamily AGE}
তিনটি আলাদা Variable।

\begin{rememberbox}

\textbf{\color{deepOrange}মনে রাখবে}

\begin{center}

\textenglish{\ttfamily name}
\(\neq\)
\textenglish{\ttfamily Name}
\(\neq\)
\textenglish{\ttfamily NAME}

\end{center}

\end{rememberbox}

%================================================
% Valid Invalid Table
%================================================

\booksubheading{Valid ও Invalid Variable Name}

\begin{tabularx}
{\linewidth}
{
|>{\RaggedRight\arraybackslash}p{0.26\linewidth}
|>{\RaggedRight\arraybackslash}p{0.20\linewidth}
|Y|
}

\hline

\rowcolor{softBlue}

\multicolumn{1}{|c|}{\textbf{Variable Name}}
&
\multicolumn{1}{c|}{\textbf{অবস্থা}}
&
\multicolumn{1}{c|}{\textbf{কারণ}}

\\

\hline

\textenglish{\ttfamily age}
&
Valid
&
Letter দিয়ে শুরু।
\\

\hline

\textenglish{\ttfamily \_age}
&
Valid
&
Underscore দিয়ে শুরু।
\\

\hline

\textenglish{\ttfamily age2}
&
Valid
&
Digit শেষে ব্যবহার করা হয়েছে।
\\

\hline

\textenglish{\ttfamily student\_name}
&
Valid
&
Underscore ব্যবহার করা হয়েছে।
\\

\hline

\textenglish{\ttfamily 2age}
&
Invalid
&
Digit দিয়ে শুরু।
\\

\hline

\textenglish{\ttfamily student name}
&
Invalid
&
Space রয়েছে।
\\

\hline

\textenglish{\ttfamily total-marks}
&
Invalid
&
Hyphen ব্যবহার করা হয়েছে।
\\

\hline

\textenglish{\ttfamily if}
&
Invalid
&
Python Keyword।
\\

\hline

\end{tabularx}

%================================================
% Meaningful Variable
%================================================

\booksubheading{অর্থবহ Variable Name ব্যবহার}

Variable-এর নাম এমন হওয়া ভালো যাতে নাম দেখেই তার কাজ বা
Value সম্পর্কে সহজে ধারণা পাওয়া যায়।

\textbf{কম পরিষ্কার—}

\begin{codebox}

\begin{english}
\ttfamily
x = 500
\end{english}

\end{codebox}

\textbf{অর্থবহ—}

\begin{codebox}

\begin{english}
\ttfamily
total\_marks = 500
\end{english}

\end{codebox}

আরও উদাহরণ—

\begin{codebox}

\begin{english}
\ttfamily
student\_name = "Mina"\\
age = 17\\
total\_price = 800
\end{english}

\end{codebox}

\begin{rememberbox}

\textbf{\color{deepOrange}সহজভাবে মনে রাখবে}

\begin{center}

\textbf{Meaningful Variable Name}
\(\longrightarrow\)
\textbf{Program বুঝতে সহজ}

\end{center}

\end{rememberbox}

%================================================
% Multiple Assignment
%================================================

\booksubheading{Multiple Assignment}

Python-এ একই Statement-এ একাধিক Variable-এর সঙ্গে
একাধিক Value Assign করা যায়।

\begin{codebox}

\begin{english}
\ttfamily
a, b, c = 10, 20, 30
\end{english}

\end{codebox}

এখানে—

\begin{center}

\textenglish{\ttfamily a = 10}
\qquad
\textenglish{\ttfamily b = 20}
\qquad
\textenglish{\ttfamily c = 30}

\end{center}

%================================================
% Same Value
%================================================

\booksubheading{একই Value একাধিক Variable-এ Assign}

একাধিক Variable-এর সঙ্গে একই Value Assign করাও সম্ভব।

\begin{codebox}

\begin{english}
\ttfamily
x = y = z = 100
\end{english}

\end{codebox}

এর ফলে—

\begin{center}

\textenglish{\ttfamily x = 100}
\qquad
\textenglish{\ttfamily y = 100}
\qquad
\textenglish{\ttfamily z = 100}

\end{center}

%================================================
% Swap
%================================================

\booksubheading{Variable-এর Value অদলবদল}

Python-এ দুটি Variable-এর Value সহজে Swap করা যায়।

\begin{codebox}

\begin{english}
\ttfamily
a = 10\\
b = 20\\[1mm]
a, b = b, a
\end{english}

\end{codebox}

এরপর—

\begin{center}

\textenglish{\ttfamily a = 20}
\qquad
\textenglish{\ttfamily b = 10}

\end{center}

\begin{rememberbox}

\textbf{\color{deepOrange}Python Shortcut}

দুটি Variable-এর Value Swap—

\begin{center}

\textenglish{\ttfamily a, b = b, a}

\end{center}

\end{rememberbox}

%================================================
% Dynamically Typed
%================================================

\booksubheading{Python Dynamically Typed}

Python-এ Variable-এর Data Type আগে আলাদাভাবে Declare
করতে হয় না।

উদাহরণ—

\begin{codebox}

\begin{english}
\ttfamily
x = 10
\end{english}

\end{codebox}

পরে একই Variable-এর সঙ্গে অন্য ধরনের Value-ও Assign করা যায়—

\begin{codebox}

\begin{english}
\ttfamily
x = "Hello"
\end{english}

\end{codebox}

অর্থাৎ Variable-এর সঙ্গে যুক্ত Value অনুযায়ী তার Data Type
পরিবর্তিত হতে পারে।

\begin{rememberbox}

এই বিষয়টি
\textbf{5.34 Python Data Type}
অংশে বিস্তারিত আলোচনা করা হবে।

\end{rememberbox}

%================================================
% Delete Variable
%================================================

\booksubheading{Variable Delete করা}

প্রয়োজনে
\textenglish{\ttfamily del}
Keyword ব্যবহার করে কোনো Variable-এর নামের binding
মুছে দেওয়া যায়।

উদাহরণ—

\begin{codebox}

\begin{english}
\ttfamily
x = 10\\
del x
\end{english}

\end{codebox}

এরপর
\textenglish{\ttfamily x}
নামটি ব্যবহার করতে গেলে Error হতে পারে।

\begin{rememberbox}

\textbf{\color{deepOrange}MCQ তথ্য}

\begin{center}

\textenglish{\ttfamily del}
\(\longrightarrow\)
\textbf{Variable-এর binding মুছে দিতে ব্যবহৃত হয়}

\end{center}

\end{rememberbox}

%================================================
% snake_case
%================================================

\booksubheading{Variable Name লেখার ভালো পদ্ধতি}

একাধিক শব্দের Variable Name লেখার ক্ষেত্রে Python-এ সাধারণভাবে
\textbf{snake\_case} পদ্ধতি ব্যবহার করা সুবিধাজনক।

উদাহরণ—

\begin{codebox}

\begin{english}
\ttfamily
student\_name = "Rahim"\\
total\_marks = 450\\
average\_score = 75.5
\end{english}

\end{codebox}

\begin{center}

\textbf{শব্দ}
+
\textenglish{\ttfamily \_}
+
\textbf{শব্দ}
\(\longrightarrow\)
\textbf{snake\_case}

\end{center}

%================================================
% Naming Rule Summary
%================================================

\booksubheading{Variable Naming Rules এক নজরে}

\begin{tabularx}
{\linewidth}
{
|>{\RaggedRight\arraybackslash}p{0.49\linewidth}
|Y|
}

\hline

\rowcolor{softBlue}

\multicolumn{1}{|c|}{\textbf{নিয়ম}}
&
\multicolumn{1}{c|}{\textbf{উদাহরণ}}

\\

\hline

Letter দিয়ে শুরু করা যায়
&
\textenglish{\ttfamily name}
\\

\hline

Underscore দিয়ে শুরু করা যায়
&
\textenglish{\ttfamily \_name}
\\

\hline

Digit পরে ব্যবহার করা যায়
&
\textenglish{\ttfamily name2}
\\

\hline

Digit দিয়ে শুরু করা যায় না
&
\textenglish{\ttfamily 2name} \; $\times$
\\

\hline

Space ব্যবহার করা যায় না
&
\textenglish{\ttfamily student name} \; $\times$
\\

\hline

Underscore ব্যবহার করা যায়
&
\textenglish{\ttfamily student\_name}
\\

\hline

Keyword ব্যবহার করা যায় না
&
\textenglish{\ttfamily if} \; $\times$
\\

\hline

Case Sensitive
&
\textenglish{\ttfamily age}
ও
\textenglish{\ttfamily Age}
আলাদা
\\

\hline

\end{tabularx}

%================================================
% Small Program
%================================================

\booksubheading{একটি ছোট Program}

\begin{examplebox}{Variable ব্যবহার}

\begin{codebox}

\begin{english}
\ttfamily
student\_name = "Mina"\\
age = 17\\
marks = 85\\[1mm]
print(student\_name)\\
print(age)\\
print(marks)
\end{english}

\end{codebox}

\textbf{Output:}

\begin{codebox}

\begin{english}
\ttfamily
Mina\\
17\\
85
\end{english}

\end{codebox}

এখানে
\textenglish{\ttfamily student\_name},
\textenglish{\ttfamily age}
ও
\textenglish{\ttfamily marks}
তিনটি Variable।

\end{examplebox}

%================================================
% Common Mistakes
%================================================

\chFiveInput{topics/33_python_variables/python_extensions.tex}

\booksubheading{সাধারণ ভুল}

\begin{compacttwocolpoints}

\item Variable Name Digit দিয়ে শুরু করা।

\item Variable Name-এর মধ্যে Space ব্যবহার করা।

\item
\textenglish{\ttfamily -},
\textenglish{\ttfamily @},
\textenglish{\ttfamily \#}
ইত্যাদি Special Character ব্যবহার করা।

\item Python Keyword-কে Variable Name হিসেবে ব্যবহার করা।

\item
\textenglish{\ttfamily age}
ও
\textenglish{\ttfamily Age}
-কে একই Variable মনে করা।

\item
\textenglish{\ttfamily =}
-কে Mathematical Equality মনে করা।

\item Value Assign করার আগে Variable ব্যবহার করার চেষ্টা করা।

\item অর্থহীন বা অস্পষ্ট Variable Name ব্যবহার করা।

\end{compacttwocolpoints}

%================================================
% Final Reminder
%================================================

\begin{rememberbox}

\textbf{\color{deepOrange}মনে রাখবে}

\begin{itemize}

\item
Python-এ Value Assign করলেই Variable তৈরি হয়।

\item
\textenglish{\ttfamily =}
হলো Assignment Operator।

\item
Variable Name Letter বা
\textenglish{\ttfamily \_}
দিয়ে শুরু হতে পারে।

\item
Variable Name Digit দিয়ে শুরু হতে পারে না।

\item
Variable Name-এর মধ্যে Space ব্যবহার করা যায় না।

\item
অধিকাংশ Special Character ব্যবহার করা যায় না।

\item
Python Keyword Variable Name হতে পারে না।

\item
Python Variable Name
\textbf{Case Sensitive}।

\item
অর্থবহ Variable Name ব্যবহার করা ভালো।

\item
একই Statement-এ একাধিক Variable-এ Value Assign করা যায়।

\item
\textenglish{\ttfamily a, b = b, a}
দিয়ে দুটি Variable-এর Value সহজে Swap করা যায়।

\item
Python
\textbf{Dynamically Typed}।

\item
দ্রুত মনে রাখবে—

\begin{center}

\textbf{Variable}
\(\;=\;\)
\textbf{Value}

\medskip

\textbf{Letter / \textenglish{\ttfamily \_}}
\(\longrightarrow\)
\textbf{শুরু করা যায়}

\qquad

\textbf{Digit}
\(\longrightarrow\)
\textbf{শুরুতে নয়}

\medskip

\textbf{Space / Special Character}
\(\longrightarrow\)
\textbf{নয়}

\qquad

\textbf{Keyword}
\(\longrightarrow\)
\textbf{নয়}

\medskip

\textenglish{\ttfamily name}
\(\neq\)
\textenglish{\ttfamily Name}

\end{center}

\end{itemize}

\end{rememberbox}

\small

%================================================
% 5.34 Python Data Type
%================================================

